%!TEX root = ../thesis.tex
% System overview for \cref{sec:impl-architecture}: the whole stack from the agent
% down to the guest binary, the two trait seams it is cut at, and the environment
% loop of \cref{sec:rl-mdp} drawn vertically through all three layers. The MOGI
% internals are one box here and are expanded in \cref{fig:platform-architecture}.
%
% Layout rules (kept deliberately consistent with figures/platform-architecture.tex,
% so the two figures read as one system at two zoom levels):
%  * A pure layered stack, agent at the top and guest binary at the bottom. The stack
%    is 106mm wide and the four component boxes in a layer are 24mm on 2mm gaps ---
%    the same numbers as the platform figure. Keep them if boxes are added: they are
%    what makes the two figures line up.
%  * Rows are chained with `below=1.5mm of', and everything else is placed relative to
%    a node border (yshift off .north, xshift off .south). Do NOT reintroduce absolute
%    coordinates: an edited label changes a row's height, and with absolute y values
%    every row below it has to be re-derived by hand.
%  * DEVIATION from the platform figure's adjacency-only rule: this figure DOES draw
%    arrows, because its subject is a flow --- one action down, one observation and
%    reward back up --- rather than a component inventory. The two channels (descent
%    on the left at -47.5mm, ascent on the right at +47.5mm) are the same circuit
%    crossing both seams, so the eye reads them as one loop. Keep them in those two
%    columns; a third arrow anywhere breaks the reading.
%  * There is deliberately NO edge between \code{efficientzero} and MOGI. That
%    \code{mugiwara} is the only layer touching both sides is a claim of the chapter,
%    not an omission.
%  * Fill encodes provenance: blue = written for this thesis, grey = existing
%    framework, white = the program under test. The caption says so; do not add a
%    legend row.
%  * The seams use `text width' rather than `minimum width'. This is not cosmetic:
%    minimum width does not clip, so an earlier draft's one-line trait list grew the
%    node past the stack and printed on top of the flow labels. The fixed 58mm is what
%    keeps the 20mm label slots on either side free.
%  * Seam lines are broken BY HAND. Automatic wrapping hyphenates identifiers --- an
%    earlier draft printed `ExitRea-son' --- so keep every seam line under ~58mm and
%    re-measure after editing one. A 24mm box wraps at ~14 characters per \scriptsize
%    line, a flow label at ~11.
%  * The phantom \path on the left balances the replication bracket on the right.
%    Without it the bounding box is asymmetric and the stack no longer sits on the
%    same page centre as \cref{fig:platform-architecture}.
\begin{figure}[htbp]
  \centering
  \begin{tikzpicture}[
      >={Stealth[round]},
      font=\small,
      layer/.style = {draw, rounded corners=2pt, fill=blue!8, align=center,
                      minimum width=106mm, minimum height=26mm},
      svc/.style   = {draw, rounded corners=1.5pt, fill=white, align=center,
                      font=\scriptsize, minimum width=24mm, minimum height=13mm},
      seam/.style  = {draw, fill=black!8, align=center, font=\scriptsize,
                      text width=58mm, inner sep=3pt},
      mogi/.style  = {draw, rounded corners=2pt, fill=black!5, align=center,
                      font=\scriptsize, minimum width=106mm, minimum height=12mm},
      guest/.style = {draw, fill=white, align=center, font=\scriptsize,
                      minimum width=106mm, minimum height=9mm},
      flow/.style  = {->, thick, black!70},
      lbl/.style   = {align=center, font=\scriptsize, text width=15mm, inner sep=1pt},
    ]
    % --- the rows, chained top to bottom -------------------------------------
    \node[layer] (ez) at (0,0) {};

    \node[seam, below=1.5mm of ez] (abs)
      {\code{efficientzero::abstraction}\\
       \code{Environment}, \code{Action},\\
       \code{Observation}, \code{StepResult}\\[1pt]
       {\color{black!55}\itshape discrete $\cdot$ deterministic $\cdot$ episodic}};

    \node[layer, below=1.5mm of abs] (mug) {};

    \node[seam, below=1.5mm of mug] (emu)
      {\code{Emulator}: \code{run} under a budget $\to$ \code{ExitReason}\\
       stdin and socket buffers; block hooks;\\
       snapshot capture and restore; fixed seeds};

    \node[mogi, below=1.5mm of emu] (mogi)
      {MOGI --- Linux user-space system emulation on Unicorn/QEMU\\
       \cref{fig:platform-architecture}};

    \node[guest, below=1mm of mogi] (guest)
      {guest: an unmodified x86-64 Linux user-space binary};

    % --- the agent: reusable, knows nothing about fuzzing ---------------------
    \node[font=\small] at ([yshift=-3.5mm]ez.north)
      {\code{efficientzero} --- the RL agent};
    \node[svc] at ([xshift=-39mm,yshift=-16.5mm]ez.north)
      {\code{driver}\\training, evaluation,\\checkpoints};
    \node[svc] at ([xshift=-13mm,yshift=-16.5mm]ez.north)
      {\code{learner}\\workers,\\replay buffer};
    \node[svc] at ([xshift= 13mm,yshift=-16.5mm]ez.north)
      {\code{mcts}\\planning using the\\learned model};
    \node[svc] at ([xshift= 39mm,yshift=-16.5mm]ez.north)
      {\code{model}\\representation,\\dynamics, policy,\\value, value prefix};

    % --- the fuzzer: implements the seam above against the seam below ---------
    \node[font=\small] at ([yshift=-3.5mm]mug.north)
      {\code{mugiwara} --- the fuzzing environment};
    \node[svc] at ([xshift=-39mm,yshift=-16.5mm]mug.north)
      {targets\\one \code{Environment}\\per program};
    \node[svc] at ([xshift=-13mm,yshift=-16.5mm]mug.north)
      {\code{corpus}\\seeds, mutation,\\action packing};
    \node[svc] at ([xshift= 13mm,yshift=-16.5mm]mug.north)
      {\code{instrument}\\block/edge hits,\\coverage buckets};
    \node[svc] at ([xshift= 39mm,yshift=-16.5mm]mug.north)
      {\code{experiments}\\probes, campaign\\persistence};

    % --- the circuit: one action down the left, observation and reward up the
    %     right. Both seams are crossed by the same two columns. ---------------
    \draw[flow] ([xshift=-47.5mm]ez.south)  -- ([xshift=-47.5mm]mug.north);
    \draw[flow] ([xshift= 47.5mm]mug.north) -- ([xshift= 47.5mm]ez.south);
    \node[lbl, left =1.5mm of abs] {action\\$a_t$};
    \node[lbl, right=1.5mm of abs] {$o_t$, $r_t$,\\done};

    \draw[flow] ([xshift=-47.5mm]mug.south)  -- ([xshift=-47.5mm]mogi.north);
    \draw[flow] ([xshift= 47.5mm]mogi.north) -- ([xshift= 47.5mm]mug.south);
    \node[lbl, left =1.5mm of emu] {input bytes,\\run, restore};
    \node[lbl, right=1.5mm of emu] {block hits,\\exit reason};

    % --- replication: one environment, emulator and guest per self-play worker
    \draw[black!55, thin]
      ([xshift=1.5mm]mug.north east) -- ([xshift=3.5mm]mug.north east) --
      ([xshift=3.5mm]guest.south east) -- ([xshift=1.5mm]guest.south east);
    \node[lbl, text width=14mm, align=left, anchor=west]
      at ([xshift=4.5mm]$(mug.north east)!0.5!(guest.south east)$)
      {one per self-play worker};

    % --- keeps the stack on the page centre despite the bracket --------------
    \path (-7.15,0);
  \end{tikzpicture}
  \caption[System architecture]{%
    The main system architecture and the fuzzing loop of~\cref{sec:rl-mdp}.
    \code{mugiwara} wraps each PUT into an environment as expected by \code{efficientzero}.
    \code{efficientzero} generates an action, \code{mugiwara} converts it to a concrete
    input, forwards the input to MOGI and receives raw execution feedback like
    block hits and the exit reason.
    \code{mugiwara} then converts this raw execution feedback to a structured observation
    and reward and forwards it to \code{efficientzero}, completing the loop.
  }
  \label{fig:system-architecture}
\end{figure}
